%%数学表达式辨别建模L的aTe5X图片、表格公式等代码模板
%%注意编译前需将编辑器的编译方式设置为xelatex
% \documentclass[withoutpreface,bwprint,utf8,a4paper]{../../../Cls/Thesis/CUMCM/CUMCM} %去掉封面与编号页
\documentclass[bwprint,a4paper]{../../../Cls/Thesis/CUMCM/CUMCM} %存储保留封面与编号页
\usepackage{../../../Styles/Thesis/CUMCM/CUMCM} %加载样式文件
\graphicspath{{figure}} %可根据个人实际情况更改图片存储位置
\rhead{题目序号}
\lhead{比赛名称}
\cfoot{\thepage}
\title{题目标题}
\tihao{A}
\membera{队员1} %第1人
\memberb{队员2} %第2人
\memberc{队员3} %第3人
\supervisor{指导老师} %指导老师
\yearinput{2025}
\monthinput{01}
\dayinput{01}
\begin{document}
    \maketitle
    \begin{abstract}
一段话。\par

另一段话。\par

一些总结
        \keywords{关键词1 \quad 关键词2 \quad 关键词3 \quad 关键词4 \quad 关键词5 }%建议4-5个
    \end{abstract}

\tableofcontents
\newpage
\section{问题重述}
\subsection{问题背景}
一个引用\cite{bib:1}。一些描述。
\subsection{相关数据}
\begin{table}[h]
    %\setlength{\abovecaptionskip}{0cm}
    %\setlength{\belowcaptionskip}{-0.2cm}
    %\large%调整整体内容的大小
    \caption{参数}
    \centering
    \begin{tabular}{cc}
        \toprule[1.5pt]
        \makebox[0.2\textwidth][c]{参数}	&  \makebox[0.35\textwidth][c]{参数值}\\
        \midrule[1pt]
        1  &  1 \\
        2 &  2 \\
        3  &  3 \\
        4  &  4\\
        5  &  5\\
        6  &  6 \\
        \bottomrule[1.5pt]
    \end{tabular}
    \label{tab:参数}
\end{table}

\subsection{需要解决的问题}
\textbf{问题一}
\par 描述。
\par
\textbf{问题二}
\par 描述。
\par
\textbf{问题三}
\par 描述。
\par
\textbf{问题四}
\par 描述。

\textbf{问题五}
\par 描述。
\par

\section{模型的假设}
    \begin{itemize}[label=]
        \item 1
        \item 2
        \item 3
        \item 4
        \item 5
        \item 6
    \end{itemize}

\section{符号的说明}
    本模型所用的各主要变量符号及意义见\cref{tab:symbol}
\begin{center}
\begin{longtable}{cc}
    \caption{主要变量符号及意义}\\
    \label{tab:symbol} \\
    \midrule
        \makebox[0.2\textwidth][c]{符号}	&  \makebox[0.6\textwidth][c]{意义}\\
        \midrule[0.5pt]
         $t$  & 时间  \\
          $r$  & 距离  \\
          $a$  & 距离 \\
          $b$  & 距离\\
          $\rho$  & 距离（单位：m）  \\
          $\omega$  &  速度 \\
          $eps$  &  误差 \\
 \bottomrule
\end{longtable}
\end{center}

\section{模型建立与求解}
\subsection{问题一的分析与求解}
\subsubsection{模型的分析}
hajkshfklajlkf4

\subsubsection{模型的建立与求解}
\paragraph{\textbf{一.aklfnlk$\theta$}} \hspace*{\fill} \par
ajflkajslf\cite{bib:2}：
\begin{equation} r(\theta) = a + b \theta
\end{equation}
\par
aklfjalkjf,$a=0$,有：
\begin{equation} r(\theta) = b \theta
\end{equation}
\par
akmfnkalf\cite{bib:2}：
\begin{equation}\label{3} x=X(\theta,b) = b\times\theta\times \cos(\theta), \quad y=Y(\theta,b) = b\times\theta \times \sin(\theta)
\end{equation}
\par 其中，$b$kaljflkasf，为$\frac{0.55}{2 \pi}$
\par
aklfklf
\newpage
\begin{thebibliography}{9}%宽度9
    \bibitem{bib:1}jasldfjk
\bibitem{bib:2}，amfnlas
\end{thebibliography}
\newpage
%附录
\begin{appendices}
    \section{求解相关结果函数}
    \begin{lstlisting}[language=matlab, basicstyle = \ttfamily,breaklines = true]

import numpy as np
import time

# Function to generate Archimedean spiral points
def archimedean_spiral(a, b, theta):
    r = a + b * theta / (2 * np.pi)
    x = r * np.cos(theta)
    y = r * np.sin(theta)
    return x, y

# 计算 getS(b, theta)
def getS(b, theta):
    sqrt_term = np.sqrt(1 + theta**2)
    log_term = np.log(theta + sqrt_term)
    b = b / (2 * np.pi)
    s_theta = 0.5 * b * (theta * np.sqrt(1 + theta**2) - np.log(theta + sqrt_term))
    return s_theta
    \end{lstlisting}


\end{appendices}

\end{document}
